\documentclass[11pt]{article}

\usepackage[a4paper,margin=2.5cm]{geometry}
\usepackage[utf8]{inputenc}
\usepackage[T1]{fontenc}
\usepackage{amsmath,amssymb,amsfonts}
\usepackage{enumitem}
\usepackage{parskip}

\title{Oefenvragen: afleidingen vanaf thermodynamica}
\author{T. Cools}
\date{}

\begin{document}
\maketitle

\section*{Thermodynamica}

\begin{enumerate}[label=\textbf{\arabic*.}, leftmargin=1.5cm]

\item Leid de toestandsfunctie enthalpie $H$ af, vertrekkend van $dU = \delta q - P_{omg}\,dV$ bij constante druk, en toon aan dat
\[
C_P = \left(\frac{\partial H}{\partial T}\right)_P.
\]

\item Leid, vertrekkend van de eerste hoofdwet bij een isotherme expansie van een ideaal gas ($U = $ cte), de uitdrukking af voor $S$ in functie van $V$:
\[
S = S^\circ + nR\ln\left(\frac{V}{V^\circ}\right).
\]

\item Leid de Gibbs vrije energie $G = H - TS$ af vanuit de tweede hoofdwet ($dS_{univ} \geq 0$), en toon aan dat
\[
dG \leq \delta w_{nuttig}
\]
bij constante $P$ en $T$.

\item Leid, vertrekkend van $G = U + PV - TS$, de uitdrukking $dG = V\,dP - S\,dT$ af, en bepaal hieruit
\[
\left(\frac{\partial G}{\partial T}\right)_P.
\]

\item Leid, met hetzelfde resultaat als in de vorige vraag, de uitdrukking af voor
\[
\left(\frac{\partial (G/T)}{\partial T}\right)_P
\]
in functie van $H$.

\item Leid voor een ideaal gas de drukafhankelijkheid van $G$ af:
\[
G(P) = G(P^\circ) + nRT\ln\left(\frac{P}{P^\circ}\right).
\]

\item Leid de Helmholtz-energie $A = U - TS$ af en toon aan dat
\[
dA = -S\,dT - P\,dV.
\]

\item Stel de Gibbs-Duhem-formule op, vertrekkend van $G$ als extensieve grootheid in $P, T, n_A, n_B$.

\item Leid de eerste Clapeyronvergelijking af ($\frac{dP}{dT}$ in functie van $\Delta H$ en $\Delta V$ van de faseovergang), vertrekkend van de voorwaarde voor fase-evenwicht en de Gibbs-Duhem-formule.

\item Leid, vertrekkend van de eerste Clapeyronvergelijking en de ideale gaswet, de Clausius-Clapeyronvergelijking af:
\[
\ln\left(\frac{P_2}{P_1}\right)
\]
in functie van $\Delta H$, $T_1$ en $T_2$.

\end{enumerate}

\section*{Chemisch evenwicht}

\begin{enumerate}[label=\textbf{\arabic*.}, leftmargin=1.5cm, resume]

\item Leid, vertrekkend van de vorderingsgraad $\xi$ en $dG = \sum_i \mu_i\,dn_i$ (P-T-controle), de uitdrukking af voor
\[
\Delta_r G = \sum_P v_P\mu_P - \sum_R v_R\mu_R.
\]

\item Leid, vertrekkend van $\mu_i = \mu_i^\circ + RT\ln(a_i)$ en de evenwichtsvoorwaarde $\Delta_r G = 0$,
de uitdrukking af voor de evenwichtsconstante $K_a^\circ$. Idem voor niet evenwichtstoestandenden.

\item Leid, voor een ideaal gasmengsel, de evenwichtsconstante $K_\chi$ af in functie van $K_P^\circ$ en de totaaldruk.

\item Leid, met het roostermodel, de entropie van een oplossing af:
\[
S = -R\left(\chi_A\ln(\chi_A) + \chi_B\ln(\chi_B)\right).
\]

\item Leid de enthalpie van een oplossing af en definieer hierbij de uitwisselingsparameter $\chi_{AB}$.

\item Leid, vertrekkend van $K_a^\circ = e^{-\Delta G^\circ/RT}$, de invloed van de temperatuur op de evenwichtsconstante af:
\[
\frac{\partial \ln(K_a^\circ)}{\partial T}.
\]

\item Leid de invloed van de druk op de evenwichtsconstante $K_\chi$ af.

\item Leid de invloed van de concentratie (via $n_A$) op de evenwichtsconstante $K_n$ af.



\end{enumerate}

\section*{Elektrochemie}

\begin{enumerate}[label=\textbf{\arabic*.}, leftmargin=1.5cm, resume]

\item Leid, vertrekkend van de arbeid van een reversibel elektrochemisch proces, de uitdrukking af voor
\[
\Delta G = -n\cdot F\cdot E_{cel}.
\]

\item Leid, vertrekkend van $\Delta G = \Delta G^0 + RT\ln(Q_A)$, de wet van Nernst af:
\[
E_{cel} = E_{cel}^0 - \frac{RT}{nF}\ln(Q_A).
\]

\end{enumerate}

\end{document}